\documentclass[11pt]{article}
\title{Temporal Signal–State Coupling in Xenopus laevis Tail Regeneration: Identifiability Limits and a Testable Receptivity-Window Hypothesis}
\date{}
\author{SIGMA AI-assisted research synthesis}
\usepackage[T1]{fontenc}
\usepackage[utf8]{inputenc}
\usepackage[margin=1in]{geometry}
\usepackage{amsmath,amssymb,array,booktabs,longtable,pdflscape}
\usepackage{microtype}
\usepackage{titlesec}
\usepackage{tikz}
\usetikzlibrary{fit}
\emergencystretch=8em
\raggedbottom
\linespread{1.07}
\usepackage{xcolor}
\definecolor{blue}{HTML}{24485E}
\titleformat{\section}{\Large\bfseries\sffamily\color{blue}}{\thesection}{0.8em}{}
\titleformat{\subsection}{\large\bfseries\sffamily}{\thesubsection}{0.7em}{}
\setlength{\parskip}{.6em}
\setlength{\parindent}{0pt}
\begin{document}
\maketitle
\begin{center}\fbox{\parbox{.92\linewidth}{\textbf{Needs review: TeX verification is incomplete.} Review the archived proofreading diagnostics before relying on the typesetting.}}\end{center}
\begin{center}\small\textit{AI-generated research synthesis \textperiodcentered{} SIGMA}\end{center}

\begin{abstract}
Why Xenopus laevis tadpoles temporarily lose tail-regeneration competence remains unresolved at the level of a discriminating, animal-level mechanism. We examine whether developmental changes in the coupling between early injury signaling and later regenerative epidermal activation could explain this switch. Under unlinked early and late sampling, identical stage-conditioned expression distributions can support opposite within-animal coupling slopes; population trajectories therefore cannot generally identify an individual predictive rule.

A restricted mechanistic model nevertheless yields a testable distinction: equal-exposure early and late epidermal ERK pulses produce opposite transition-drive contrasts under insufficient-signal and delayed-receptivity explanations. This prediction depends on selective delivery, validated state measurements, and separation of wound age from response latency. No measured biological effect, baseline improvement, or independent confirmation supports the proposal. The contribution is a conditional identifiability argument and a bounded experimental hypothesis, not a discovered mechanism or an explanation of later regained competence.
\end{abstract}
\section*{About this paper}
This paper was generated by SIGMA, an AI-assisted research system, from the project's recorded investigations, findings, and supporting evidence.
The analysis and conclusions reflect an automated synthesis of that research. Although SIGMA incorporates evidence tracking, critical review, and consistency checks, the resulting paper has not necessarily undergone independent verification or peer review.
Readers should evaluate its arguments, supporting evidence, and stated limitations before relying on its conclusions.
Research style: Distinguish observed biology from mechanistic hypotheses and proposed experiments. Identify organisms, developmental stages, controls, and cross-species limitations.
\section*{Evidence coverage limitation}
The evidence index omitted 3370 material saved research field fragment(s) across 10 source(s). These fragments were not reviewed for this paper. Interpretations that could depend on omitted material remain provisional; the complete frozen source snapshot and omission ledger are retained with the generation record.
\textit{Claims discussed: }C3, C5, C6, C7, C11, C12
\section{Introduction and biological context}
The biological target is the temporary loss and subsequent recovery of tail-regeneration competence in Xenopus laevis, including differences between animals at nominally similar developmental stages. The motivating comparison identifies stage 40 as regeneration competent and stage 46 as refractory, but stage membership is not an observed regeneration outcome for each animal. An explanation of the full developmental switch must also account for a separately verified later competent condition.

Transient regeneration-initiating cells, abbreviated RICs, and extracellular-matrix remodeling are treated as established prior-art exclusions in the supplied biological background, not as discoveries of this analysis. The mechanistic question is narrower: does developmental stage alter when basal epidermal cells can convert injury signaling into a regenerative state? The proposed contribution concerns this conditional timing relationship rather than the involvement of a familiar pathway.

A useful predictor must distinguish signal insufficiency, altered signal-to-state coupling, metabolic or tissue-composition gating, and analytical confounding. It must improve on stage-only, RIC/ECM-only, and known-signaling baselines using biological units with appropriate outcome linkage. Evidence from Xenopus tropicalis cannot substitute for species-specific evidence about the Xenopus laevis developmental trajectory.

\textit{Claims discussed: }C1, C2, C5, C10, C11
\section{Methods and evidence}
The analysis formalized temporal coupling as a within-stage relationship between an animal's early injury-response measurement and its later regenerative epidermal measurement. Explicit observational-equivalence constructions were examined to determine whether separate early and late samples identify that relationship. A second algebraic argument examined whether assigning stage-average early responses to late samples creates an estimable stage-by-response interaction.

A restricted signal\textendash{}receptivity model compared insufficient early signaling with delayed epidermal receptivity. Its equal-exposure pulse contrast was derived at the level of cumulative transition drive, and the assumptions needed to extend that contrast to state induction or regrowth were separated from the derivation. Saved objections to redox persistence, spatial alignment, and fixed-time molecular readouts were incorporated as design constraints; no expression reanalysis, baseline fitting, or experimental testing was performed.

\textit{Claims discussed: }C3, C4, C6, C7, C8, C9, C12
\section{The mechanistic question beyond known regenerative pathways}
The mechanistic question is whether developmental loss of tail-regeneration competence in Xenopus laevis reflects a change in when injured epidermal cells can respond to signaling, rather than simply whether a familiar regenerative pathway is activated. The motivating comparison is between stage-40 tadpoles described as regeneration-competent and stage-46 tadpoles described as refractory. Because competence varies between animals, the target is not merely a molecular distinction between stage classes, but a mechanism that could account for differences in regenerative response within that developmental contrast.

These stage descriptions motivate the hypothesis; they do not establish a timing-dependent mechanism or an individual predictor. [C5, C11]

The supplied biological background places regeneration-initiating cells (RICs) and extracellular-matrix (ECM) remodeling outside the novelty claim. It attributes transient RIC formation from basal epidermal cells and remodeling defects associated with regeneration failure to the 2024 Genome Biology study identified by DOI 10.1186/s13059-024-03396-3, and treats reactive oxygen species and inflammatory and developmental signaling as established regeneration-related processes. Accordingly, neither detecting these cells nor implicating those pathways would answer the mechanistic question posed here.

The original publications and subsequent literature have not been examined sufficiently to establish whether the specific timing interaction proposed below is previously unreported; this background therefore defines an exclusion from novelty, not evidence of priority for the present hypothesis. [C1]

The candidate explanation concerns temporal overlap between an injury signal and basal-epidermal receptivity to entering the transient regenerative state. Under a restricted insufficient-signal explanation, stage-46 epidermal cells remain receptive early after injury but receive inadequate early signaling. Under the alternative delayed-receptivity explanation, those cells may remain inducible, yet become receptive only after endogenous signaling has declined. Both could produce failure to induce the regenerative epidermal state, so that shared endpoint alone would not distinguish them. Neither temporal arrangement is an observed biological result in the supplied evidence. [C5]

This distinction shifts the proposed contribution from pathway membership to a conditional relationship: the effectiveness of basal-epidermal ERK activity may depend jointly on developmental stage and time after injury. The relevant question is therefore whether comparable signaling delivered at different wound ages has different effects on state entry, not simply whether more signaling accompanies more regeneration. Matching signal exposure is conceptually important because an apparent timing advantage caused by a larger signal would not separate the restricted amplitude and receptivity explanations. This formulation identifies a bounded mechanistic comparison without assuming that ERK timing explains regeneration failure. [C5]

Later recovery of competence is a separate biological requirement, not a consequence that can be inferred from the stage-40 versus stage-46 contrast. No verified later-competent X. laevis condition in the supplied evidence supports the proposed receptivity mechanism, and stage patterns or nutritional findings attributed to X. tropicalis cannot supply that missing comparison. The argument is therefore narrowed from explaining the entire loss-and-recovery trajectory to distinguishing insufficient early signaling from mistimed epidermal receptivity during the motivating loss-of-competence contrast.

Explaining restoration would additionally require showing how the relevant signal\textendash{}state relationship changes in a confirmed later-competent X. laevis condition. [C11]

\textit{Claims discussed: }C1, C5, C11
\section{What unlinked transcriptomic snapshots can identify}
Transcriptomic snapshots can characterize stage-dependent population responses without identifying whether an individual animal\textquoteright{}s early injury response predicts its later regenerative epidermal activation. This distinction determines what the proposed Xenopus laevis resource can contribute to a temporal-coupling explanation of competence. GSE245320 is described as encompassing bulk RNA-seq series GSE245317/GSE245318, single-cell series GSE245312, and spatial series GSE245313; the motivating bulk comparison is described as stages 40 and 46 sampled at $0$, $0.5$, and $6$ hours and $3$ days after amputation. These descriptions identify candidate contrasts, not an established sampling design. [C2, C3]

The usable design remains unresolved because an inspected sample manifest and expression files are not available in the supplied evidence. Essential unknowns include animal versus pool identity, pool size, biological and technical replication, exact stage and time coverage, tissue boundaries, expression units, clutch and batch structure, file availability, early-to-late correspondence, and regeneration-outcome labels. A library-level manifest must establish these features before either population contrasts or animal-level models are selected: technical libraries from one specimen do not constitute additional biological replicates, and independently collected pools cannot be treated as individual animals. [C2]

The modalities should consequently be regarded as complementary measurements from a shared study, not independent biological confirmation merely because they have different accession numbers. The single-cell resource is described as predominantly characterizing regenerating animals rather than providing matched snapshots across both developmental conditions. Such measurements may help define candidate epidermal states, but they cannot supply missing refractory observations or recover early-to-late animal correspondence without additional validated information. Likewise, linking measurements at the pool level would support a pool-level estimand, not automatically an animal-level predictor. [C2, C3]

To specify the individual-level target, let $E_i$ denote animal $i$\textquoteright{}s early wound-response measurement, $R_i$ its later basal-epidermal regeneration-initiating-cell/extracellular-matrix (RIC/ECM) transition measurement, and $S_i$ its developmental stage. With finite relevant moments and positive within-stage early-response variance, define the coupling slope as $\beta_s=\operatorname{Cov}(E,R\mid S=s)/\operatorname{Var}(E\mid S=s)$. A stage-dependent coupling claim concerns a difference such as $\Delta\beta=\beta_{46}-\beta_{40}$: whether the within-stage relationship between early response and later activation changes with development. This is an association estimand, not by itself a causal effect of injury signaling. [C3]

If early and late samples are independent, unlinked biological units, their observations identify the marginal distributions $P(E\mid S=s)$ and $P(R\mid S=s)$ rather than the joint distribution $P(E,R\mid S=s)$. The missing information is which later response belongs to the animal with a particular early response. Matching cohorts by stage, batch, or clutch does not create that individual correspondence. The following limitation therefore applies when only these marginals are available and no additional justified restrictions identify their joint distribution; it is not an established description of the named accessions. [C2, C3]

An exact counterexample shows why the missing correspondence matters even if both marginal distributions are known perfectly. Within a stage, let $U\sim\operatorname{Uniform}[0,1]$: one construction assigns $E=U$ and $R=U$, whereas another assigns $E=U$ and $R=1-U$. Both produce identical uniform early and late marginals, but $\operatorname{Var}(E)=1/12$ and the respective covariances are $1/12$ and $-1/12$, yielding slopes $+1$ and $-1$. Thus the observable distributions under independent, unlinked sampling are identical despite opposite signs of individual coupling; these are constructed mathematical values, not measured tadpole expression. [C3]

The same construction makes the stage interaction nonidentifiable. Assigning the positive-coupling construction to stage 40 and the negative-coupling construction to stage 46 gives $\Delta\beta=-2$; reversing the assignments gives $\Delta\beta=+2$, without changing any stage-conditioned early or late marginal. Adding stage-specific means or positive scale factors does not remove the underlying ambiguity about correspondence. Consequently, a reproducible stage-by-time expression contrast can coexist with opposing within-animal coupling rules, and additional unlinked samples or time points do not generally resolve that ambiguity without further identifying information. [C3]

Replacing missing individual early measurements with stage-specific early means introduces a separate, algebraic failure. Code the two stages as $S\in\{0,1\}$ and assign every late sample $E=e_0+(e_1-e_0)S$, where $e_0$ and $e_1$ are the respective early means. Since $S^2=S$, the interaction satisfies $SE=e_1S$, so the intercept, stage, assigned early response, and stage-by-early-response columns span at most two dimensions. More late biological replicates may sharpen the stage means, but cannot identify a separate early-response coefficient or interaction from this rank-deficient design. [C4]

This boundary preserves a useful role for cross-sectional transcriptomics while changing the interpretation of its contribution. Where replication and stage\textendash{}batch separation support them, differential-expression and stage-by-time analyses can characterize population trajectories; correlations among group means remain population-level relationships rather than estimates of $\beta_s$. Actual longitudinal links, informative interventions, or justified structural restrictions could change identifiability.

The central argument therefore turns on distinguishing a developmental difference in average activation from a developmental change in signal\textendash{}state coupling: the former can motivate the mechanism, but cannot substitute for evidence of the latter. [C2, C3, C4]

\textit{Claims discussed: }C2, C3, C4
\section{A restricted signal–receptivity model}
A restricted signal\textendash{}receptivity model makes the timing hypothesis discriminating: stage-46 failure could reflect either insufficient early epidermal signaling or a receptive state that becomes available only after that signaling has declined. For Xenopus laevis tadpoles, the proposed distinction concerns basal-epidermal entry into a transient regenerative state, rather than the presence of a regeneration-associated pathway alone. Delayed receptivity is an unconfirmed candidate mechanism, not an observed property of stage-46 epidermis [C5].

Let $E_s(t)$ denote basal-epidermal ERK activity at developmental stage $s$ and time $t$ after amputation, and let $C_s(t)$ denote receptivity to entering the transient regenerative epidermal state. Define cumulative transition drive as $J_s=k\int E_s(t)C_s(t)\,dt$, with $k>0$, over a common post-amputation interval. This definition makes the temporal overlap between activity and receptivity consequential: activity contributes to the modeled transition only when receptivity is present. Here, $C_s(t)$ is a latent mechanistic quantity whose operational measurement remains unresolved; the equation stipulates a candidate relationship rather than estimating one from biological data [C5, C6].

In the restricted insufficient-signal model, receptivity is concentrated in an early window at both stage 40 and stage 46, but stage-46 ERK activity is insufficient during that window. In the restricted delayed-receptivity model, stage 40 retains early receptivity, whereas stage-46 receptivity occurs in a later window, potentially after endogenous ERK activity has declined. Both models can therefore yield insufficient transition drive and failure to induce the regenerative epidermal state. The broad description \textquotedblleft{}failure to induce RICs\textquotedblright{} encompasses both explanations and does not distinguish deficient signal amplitude from deficient temporal overlap [C5].

To derive a distinguishing intervention contrast, consider two nonnegative additions to stage-46 ERK activity, $p_{\mathrm{early}}(t)$ and $p_{\mathrm{late}}(t)$, supported within their respective, nonoverlapping windows. Require matched added exposure, $\int p_{\mathrm{early}}(t)\,dt=\int p_{\mathrm{late}}(t)\,dt=Q$, where $Q>0$. These pulses represent added cellular ERK activity, not merely equal durations of an external stimulus. The analytical comparison holds the untreated activity trajectory and receptivity profile fixed between the two timing conditions [C6].

Writing the common untreated trajectory as $E_{46}^{0}(t)$ gives $J_{46,\mathrm{late}}=k\int [E_{46}^{0}(t)+p_{\mathrm{late}}(t)]C_{46}(t)\,dt$ and $J_{46,\mathrm{early}}=k\int [E_{46}^{0}(t)+p_{\mathrm{early}}(t)]C_{46}(t)\,dt$. Subtraction explicitly cancels the shared untreated contribution: $\Delta J=J_{46,\mathrm{late}}-J_{46,\mathrm{early}}=k\int E_{46}^{0}(t)C_{46}(t)\,dt-k\int E_{46}^{0}(t)C_{46}(t)\,dt+k\int C_{46}(t)[p_{\mathrm{late}}(t)-p_{\mathrm{early}}(t)]\,dt$. Thus $\Delta J=k\int C_{46}(t)[p_{\mathrm{late}}(t)-p_{\mathrm{early}}(t)]\,dt$: the remaining contrast depends on where the matched additions overlap receptivity [C6].

For the idealized insufficient-signal model, let stage-46 receptivity equal a constant $c>0$ in the early window and zero in the late window. The late pulse then contributes no additional transition drive, whereas the early pulse contributes $kcQ$, giving $\Delta J=-kcQ$. For the idealized delayed-receptivity model, reversing these receptive supports gives $\Delta J=+kcQ$. These opposite signs are exact consequences of the restricted models at equal added exposure, not measured effect sizes [C6].

The sign result depends on those temporal-support assumptions, not on the labels assigned to the mechanisms. For a general receptive profile, the contrast remains the pulse-weighted difference in receptivity, and neither sign is guaranteed merely by describing a defect as an amplitude failure or a delay. Receptivity absent from both tested windows would yield no added drive under either pulse. Consequently, the proposed comparison discriminates between the stipulated receptive profiles, not between all possible injury-signaling defects and all possible RIC or extracellular-matrix defects [C5, C6].

A transition-drive contrast also does not automatically imply an equivalent regrowth contrast. Transferring its sign to regrowth requires a monotone, nonsaturated downstream response without another dominant block; saturation or an independent downstream limitation can erase the predicted difference. The model therefore changes the mechanistic question from whether epidermal signaling can promote regeneration to whether the same added activity has opposite timing preferences under two specified explanations. That conditional timing preference\textemdash{}not rescue by ERK supplementation alone\textemdash{}is the discriminating prediction [C6].

\begin{quote}
\textbf{Diagram (readable relationship list).} The restricted models predict opposite effects of pulse timing on cumulative transition drive. Early and late added ERK pulses have the same exposure integral; their effect depends on whether basal-epidermal receptivity occurs early or late. These are hypothetical model contrasts, not observed treatment effects, and they do not by themselves determine regeneration.
\begin{description}
\item[\textbf{Node 1.}] Early and late epidermal ERK pulses with equal integral $Q>0$
\item[\textbf{Node 2.}] Restricted signal-insufficiency model: stage-46 receptivity is early
\item[\textbf{Node 3.}] Restricted delay model: stage-46 receptivity is late
\item[\textbf{Node 4.}] Late-minus-early drive: $\Delta J=-kcQ$
\item[\textbf{Node 5.}] Late-minus-early drive: $\Delta J=+kcQ$
\item[\textbf{Node 6.}] State induction and regrowth require validated readouts, selective delivery, and
\item[\textbf{Edge 1.}] Node 1 points to node 2: Evaluate overlap with the early receptive window
\item[\textbf{Edge 2.}] Node 1 points to node 3: Evaluate overlap with the late receptive window
\item[\textbf{Edge 3.}] Node 2 points to node 4: Only the early pulse overlaps receptivity
\item[\textbf{Edge 4.}] Node 3 points to node 5: Only the late pulse overlaps receptivity
\item[\textbf{Edge 5.}] Node 4 points to node 6: Translation to a phenotype is conditional
\item[\textbf{Edge 6.}] Node 5 points to node 6: Translation to a phenotype is conditional
\item[\textbf{Boundary 1.}] Analytical predictions under restricted, unconfirmed mechanisms: node 1, node 2, node 3, node 4, node 5
\item[\textbf{Boundary 2.}] Unresolved experimental and downstream assumptions: node 6
\end{description}
\par\smallskip{\footnotesize Evidence: [Research session 21, artifact 4], [Research session 21, artifact 5]; claims: C5, C6, C7, C12}
\end{quote}
\textit{Claims discussed: }C5, C6
\section{Discriminating timing from exposure, kinetics, and spatial organization}
A pulse-timing effect would support wound-age-dependent receptivity only if it could be separated from unequal response latency, integrated exposure, and spatial organization. For the proposed Xenopus laevis stage-40 and stage-46 comparison, these are distinct identification requirements: equal measured exposure does not equalize time available for induction or decay, and equal tissue-level activity does not establish where signaling and remodeling occur. The receptivity-window hypothesis therefore requires a temporal response pattern that competing kinetic, exposure-based, and spatial explanations cannot reproduce, rather than merely a larger molecular score after a later pulse. [C7\textendash{}C9]

The fixed six-hour molecular endpoint fails the first requirement because changing pulse time also changes elapsed time since treatment. If a pulse is anchored at wound age $p$ and the response is measured at wound age $t$, its post-pulse latency is \$\textbackslash{}ell=t-p\$; holding $t$ fixed necessarily gives early and late pulses different latencies. A transient response could consequently be lower after an early pulse because it has already decayed, even when early and late pulses have identical efficacy at matched post-pulse times. Stage-dependent induction or decay kinetics could also produce a stage-by-pulse-timing interaction without any developmental shift in the receptive window. [C7]

The proposed experiment therefore still needs a validated sampling schedule that distinguishes wound age from elapsed post-pulse time. Matched-latency measurements should compare early and late interventions after equal response intervals, while vehicle-adjusted responses should use controls measured at the corresponding wound ages. Measurements spanning induction and subsequent persistence are needed to test a kinetic alternative in which pulse efficacy is unchanged but the observed score rises or decays differently between stages. A set of sampling times that merely brackets the intervention windows does not, by itself, establish this separation; the schedule must make the relevant matched-latency and vehicle comparisons possible. [C7]

These timing comparisons also depend on intervention and measurement capabilities that remain to be established in basal epidermis. For the proposed ERK experiment, epidermis-restricted delivery, reversible activation or inhibition, noninjurious treatment, and comparable achieved activity trajectories are feasibility obligations, not demonstrated properties of an available protocol. Matching illumination duration or nominal reagent dose would not substitute for matching measured local signaling exposure, and systemic inhibition would not establish basal-epidermal necessity.

The state readout must likewise resolve induction per eligible basal-epidermal cell, alongside cell abundance and subsequent remodeling; an apparent absence of state induction can challenge the proposed mediation only if the assay covers the relevant cells and times with adequate sensitivity. [C7]

The redox-resolution alternative illustrates a separate failure to distinguish temporal organization from total exposure. Its biological proposal is that, conditional on an adequate initial wound response, persistent basal-epidermal peroxide signaling between half an hour and six hours inhibits the subsequent RIC/ECM transition. This is an untested candidate, not an observed peroxide trajectory or a demonstrated mechanism in refractory animals. The proposed perturbation would shorten the later response while preserving the initial response, but that contrast alone would also alter integrated oxidant exposure. [C8]

More importantly, the proposed persistence score is mathematically an exposure measure. Writing baseline-adjusted peroxide-reporter activity as $h(t)$, with time measured in hours, the score is $P=\frac{1}{5.5}\int_{0.5}^{6}h(t)\,dt$. Matching the measured integral over that interval necessarily fixes $P$, so an integral-matched short-tail versus long-tail comparison cannot test a persistence-specific rule expressed solely through this score. A corrected test would first need a separate temporal-shape estimand capable of varying at fixed integral, together with feasible interventions that alter that feature selectively; neither is established here.

This is a failure of the discriminating design, not evidence that the redox-resolution hypothesis is false. [C8]

Spatial injury\textendash{}ECM alignment remains another competing explanation because marginal activity measurements discard the correspondence between wound positions. Consider the supplied mathematical construction of four equally weighted, registered positions, with early injury activity $I=(1,1,0,0)$ and alternative later ECM fields $M_A=(1,1,0,0)$ or $M_B=(0,0,1,1)$. Both alternatives have the same total activity, active-position count, and mean, yet their centered spatial products, $Q_{\mathrm{raw}}=\overline{IM}-\bar I\bar M$, are respectively $1/4$ and $-1/4$.

The construction establishes that equal abundance can conceal opposite alignment; it neither demonstrates these arrangements in tadpoles nor shows that alignment affects regeneration. [C9]

Testing this spatial candidate would require registered measurements linking early basal-epidermal injury activity to later local ECM activity, rather than co-detection in an unresolved tissue mixture. The proposed half-hour injury and six-hour ECM measurements are prospective choices whose suitability, reporter specificity, and registration accuracy still require validation. A candidate alignment statistic can be written as $Q=T-\mathbb{E}_{\mathrm{null}}[T]$, where $T=\overline{IM}$ and the null expectation represents alignment attributable to specified anatomical constraints.

Its interpretation would depend on reliable cell identity, geometry, and composition measurements, so that segmentation or tissue-mixture artifacts do not masquerade as local correspondence. [C9]

The spatial null must preserve more than marginal distributions. Reassigning measurements among positions within wound-distance and anatomical strata may retain activity totals and cell counts while changing spatial autocorrelation and patch structure; shared smooth anatomical gradients could then resemble excess alignment. The comparison therefore needs explicit exchangeability assumptions, specification of which spatial features are retained, and sensitivity checks for registration error and stratum choice. Even a well-calibrated observational alignment association would not establish a causal coupling defect, because anatomical or metabolic common causes could generate both alignment and regenerative success. [C9]

These limitations narrow what the proposed temporal contrast could establish. A later-pulse advantage at a common endpoint is not yet evidence for delayed receptivity, an integral-derived score cannot distinguish persistence from burden under integral matching, and abundance controls cannot exclude spatial misalignment. The central argument consequently shifts from the sign of a single endpoint contrast to whether validated, exposure-matched perturbations produce vehicle-adjusted state responses at matched latency that cannot be explained by the specified kinetic alternative, while spatial organization remains a separately testable competitor. [C7\textendash{}C9]

\textit{Claims discussed: }C7, C8, C9
\section{Evidence needed for prediction and a developmental mechanism}
Testing temporal signal\textendash{}state coupling requires three distinct evidentiary targets: an identifiable population association, an individual regeneration predictor, and a perturbation-supported developmental mechanism. Each requires information that the preceding target does not supply; neither a stage-dependent expression trajectory nor improved prediction alone establishes how injury signaling controls regenerative state entry. The supplied evidence contains no expression reanalysis, fitted baseline, measured biological effect, or independent functional confirmation, so the requirements below define prospective tests rather than demonstrated performance [C10, C12].

For population association, the first requirement is a replicated comparison whose stage and sampling-time effects can be distinguished from batch and biological-unit differences. Before fitting a stage-by-time model, sample records must establish whether observations represent individual animals, independently collected pools, or technical measurements of the same specimen. If early and later collections are unlinked, analysis should remain restricted to identifiable population contrasts rather than treating group trajectories as individual temporal coupling. A pool-level association, where identifiable, must likewise retain its pool-level interpretation [C10].

An individual predictor requires molecular measurements genuinely available from the animal before its regeneration outcome is scored. A prospective design would therefore obtain an early wound-response measurement and a later basal-epidermal state measurement nonterminally in the same Xenopus laevis tadpole, followed by blinded regeneration scoring. Reporter specificity, kinetics, and sufficiently nonperturbative measurement would need validation before these measurements could support prediction. Matching separate destructive cohorts by stage, clutch, or treatment does not supply the missing individual measurement\textendash{}outcome links [C10].

This feasibility requirement also determines which baseline comparisons are legitimate. The preregistered comparator set should include stage-only, regeneration-initiating-cell/extracellular-matrix (RIC/ECM)-only, and known-signaling-module models, with stage-adjusted molecular comparators used to assess information beyond developmental stage. A destructive RIC assay cannot serve as an individual-level baseline for another animal followed to regeneration; that comparison requires a validated live surrogate or an explicitly narrower, feasible baseline set. Any claimed improvement must be evaluated on identical eligible test animals using measurements available at the declared prediction time [C10].

Predictive evaluation should separate development of the candidate rule from assessment on untouched biological units. All cells, images, and technical libraries from an animal must remain together, and independently collected pools must not be split into apparent animal-level replicates. Held-out clutches and batches should test performance beyond the conditions used to select modules and fit models, with normalization and scaling estimated from training data alone. Evidence of an association after adjustment is a different result from improved held-out prediction; the latter requires a prespecified performance metric and uncertainty for the difference from each relevant baseline [C10].

Composition controls are essential to the interpretation of any molecular advantage. The analysis should distinguish basal-cell abundance and total tissue signal from state induction within eligible basal epidermal cells, while accounting for prespecified wound-size, closure, viability, developmental, and batch variables. Where cell composition is inferred rather than directly measured, its uncertainty should be retained and the estimates checked against histology. Attenuation after composition adjustment would favor a tissue-mixture explanation over a cell-intrinsic coupling interpretation, but persistence after adjustment would not exclude metabolic, anatomical, or injury-related common causes [C10].

Preregistration should fix molecular modules, X. laevis homeolog handling, outcome scoring, exclusions, baseline models, and the primary comparison before confirmatory outcomes are examined. A single primary interaction or predictive comparison should be distinguished from secondary modules, sampling times, and endpoints, with multiplicity handling specified for those additional tests. Negative controls should address the same proposed source of predictive information rather than merely reproduce a stage difference. Similar candidate and control point estimates would not establish equivalence without a prespecified margin and an uncertainty rule [C10].

Precision planning must reflect between-animal and between-clutch variation, measurement error, and outcome imbalance rather than assuming that a proposed animal count is sufficient. A separate pilot should justify a biologically meaningful minimum effect, establish feasible measurement times, and inform a finite confirmatory sample size without using confirmatory outcomes to select the response module or threshold. The analysis should report clutch-aware intervals for the primary association, predictive improvement, or perturbation contrast, as appropriate. For individual prediction beyond stage classification, the evaluation also requires sufficient within-stage outcome variation; distinguishing stage labels is not itself the target [C10].

A developmental mechanism requires independent intervention evidence linking signal timing to epidermal state induction and regeneration. Newly raised X. laevis from multiple clutches would permit treatment randomization within stage and clutch, with molecular and regeneration outcomes evaluated separately and with uncertainty. Terminal molecular cohorts can test treatment effects on state induction, while separately randomized animals can test regeneration effects, but those arms cannot be joined into an individual longitudinal predictor. Even concordant effects would require controls for generic growth, toxicity, wound severity, composition, and developmental progression before being attributed to the proposed coupling mechanism [C10].

The causal distinction is also narrower than rescue alone. Early-signal perturbation should be accompanied by a downstream bypass or other discriminating control before an association is attributed to failed signal transmission rather than an inability to enter the RIC state. Sham, vehicle, wound-size, closure, viability, and composition controls are needed to interpret that contrast. A nonsignificant rescue result remains inconclusive unless manipulation and assay checks are valid and the interval excludes a prespecified meaningful benefit; conversely, rescue without detected state induction challenges the proposed mediation only if the assay covers the relevant cells and times with adequate sensitivity [C10].

Finally, the proposed stage-40 versus stage-46 contrast addresses loss of competence, not its subsequent return. Extending the mechanism to a loss-and-recovery trajectory requires an independently verified later-competent X. laevis condition, followed by tests of the proposed signal\textendash{}state relationship in that condition; stage patterns or nutritional findings from X. tropicalis cannot substitute for this comparison [C11]. The evidentiary consequence is a separation of conclusions: replicated population contrasts can characterize developmental responses, linked pre-outcome measurements can test individual prediction, and selective perturbations with independent regeneration outcomes can test the restricted mechanism.

The current analytical result therefore changes what must be measured to advance the receptivity-window hypothesis, rather than establishing that the hypothesis explains the developmental switch [C10, C12].

\textit{Claims discussed: }C10, C11, C12
\section{Findings and hypotheses}
\subsection{C1 -- Limitation}
The supplied background excludes RIC identification, ECM-remodeling involvement, and familiar injury-signaling pathways from novelty, but the original publications and subsequent literature have not been inspected sufficiently to establish the novelty of a timing interaction.
The supplied account attributes basal-epidermal RIC origins and ECM-remodeling defects to DOI 10.1186/s13059-024-03396-3. DOI 10.1016/j.ydbio.2021.01.005 and DOI 10.3390/ijms252111597 are additional reference pointers whose detailed contents remain unchecked. These pointers establish a provisional novelty boundary, not a completed literature review or priority claim.
\textit{Evidence: }[Project overview], [Research session 21, artifact 1], [Research session 21, artifact 4]
\subsection{C2 -- Unresolved}
The actual sampling structure and usable contrasts of GSE245320 and its named component series remain unverified.
The project identifies bulk series GSE245317/GSE245318, single-cell series GSE245312, and spatial series GSE245313 as components of a shared-study resource. The reported stages and collection times are starting descriptions rather than an inspected sample manifest. Animal versus pool identity, replication, units, batches, pairing, phenotype labels, spatial resolution, and file availability have not been established; the described single-cell resource is not a verified matched comparison across both conditions.
\textit{Evidence: }[Project overview], [Research session 21, artifact 5], [Research session 21, artifact 3]
\subsection{C3 -- Established}
Under unlinked sampling that identifies only stage-conditioned early and late marginal distributions, an animal-level temporal-coupling coefficient and its stage interaction are not generally identifiable.
Let $E$ be an animal's early response, $R$ its later RIC/ECM measurement, and $S$ developmental stage. With finite moments and positive within-stage early variance, define $\beta_s=\operatorname{Cov}(E,R\mid S=s)/\operatorname{Var}(E\mid S=s)$. For a constructed variable $U$ uniform on $[0,1]$, the assignments $E=U,R=U$ and $E=U,R=1-U$ give identical uniform marginals. Since $\operatorname{Var}(E)=1/12$, their covariances are respectively $1/12$ and $-1/12$, and their slopes are $+1$ and $-1$. Assigning these joint distributions differently across stages changes the coupling interaction without changing any observed marginal. This establishes a conditional analytical limitation, not its applicability to the named accessions and not a biological finding.
\textit{Evidence: }[Research session 21, artifact 2], [Research session 21, artifact 5], [Research session 21, artifact 1]
\subsection{C4 -- Established}
Replacing missing individual early measurements with stage-specific early means cannot identify a separate stage-by-early-response interaction in a two-stage regression.
Code the two stages as $S\in\{0,1\}$ and assign each late sample the value $E=e_0+(e_1-e_0)S$. The product then satisfies $SE=e_1S$. Consequently, the intercept, stage, assigned early response, and interaction columns span at most two dimensions. More late samples may improve estimates of stage means but cannot restore the missing independent predictor variation.
\textit{Evidence: }[Research session 21, artifact 2]
\subsection{C5 -- Conjecture}
A delayed basal-epidermal receptivity window is an unconfirmed candidate explanation for stage-46 regeneration failure, distinct from a restricted insufficient-early-signal model.
Let $E_s(t)$ denote basal-epidermal ERK activity and $C_s(t)$ receptivity to entering the transient regenerative epidermal state. The restricted amplitude model retains early receptivity at both stages but attributes stage-46 failure to insufficient early activity. The restricted delay model retains early receptivity at stage 40 but shifts stage-46 receptivity to a later interval, potentially after endogenous activity has declined. Neither model is observed, and the broad phrase 'failure to induce RICs' includes both rather than distinguishing them.
\textit{Evidence: }[Research session 21, artifact 4]
\subsection{C6 -- Established}
Under the restricted receptivity models and matched added ERK exposure, late-minus-early supplementation has opposite signs for cumulative transition drive.
Define cumulative drive by $J_s=k\int E_s(t)C_s(t)\,dt$, with $k>0$. Let nonnegative added pulses $p_{\mathrm{early}}$ and $p_{\mathrm{late}}$ occupy their respective windows and each have integral $Q>0$, while the untreated background and receptivity remain fixed within stage. Subtraction cancels the background and gives $\Delta J=k\int C_{46}(t)[p_{\mathrm{late}}(t)-p_{\mathrm{early}}(t)]\,dt$. If receptivity is $c>0$ only in the early window, then $\Delta J=-kcQ$; if it is $c$ only in the late window, then $\Delta J=+kcQ$. These are exact model contrasts, not measured effect sizes. General receptive profiles need not preserve either sign, and a corresponding regrowth contrast additionally requires a monotone, nonsaturated downstream response without another dominant block.
\textit{Evidence: }[Research session 21, artifact 4]
\subsection{C7 -- Limitation}
A fixed six-hour molecular endpoint cannot by itself distinguish wound-age-dependent receptivity from differences in time elapsed after early and late pulses.
Early and late interventions have different post-treatment latencies when measured at one common wound age. Transient induction, score decay, or stage-dependent response kinetics can therefore generate a timing contrast without a shifted receptive window. Matched-latency, vehicle-adjusted measurements and a schedule that separates these two time coordinates remain necessary design obligations.
\textit{Evidence: }[Research session 21, artifact 5], [Research session 21]
\subsection{C8 -- Limitation}
The proposed redox-resolution alternative has not been tested, and an exposure-matched experiment cannot distinguish persistence from total exposure when its persistence score is exactly proportional to the matched integral.
The saved biological candidate proposes that, conditional on an adequate initial wound response, persistent epidermal peroxide signaling between half an hour and six hours inhibits the later RIC/ECM transition. Its objection identifies a persistence score proportional to the same interval's exposure integral. Holding that integral fixed also fixes the score, eliminating the proposed discriminating contrast. This is a design failure, not evidence that the biological redox hypothesis is false.
\textit{Evidence: }[Research session 21, artifact 2]
\subsection{C9 -- Unresolved}
Spatial injury\textendash{}ECM alignment remains a competing hypothesis that abundance measurements alone cannot establish, and its proposed null comparison must address spatial structure.
The saved construction places identical amounts of injury and ECM activity either at the same or at different wound positions, showing that equal marginal abundance need not imply equal alignment. It does not demonstrate either arrangement in tadpoles or establish an effect on regeneration. Reassignment within anatomical strata may change autocorrelation and patch structure, so shared anatomical gradients can mimic alignment unless the null's exchangeability and retained spatial structure are specified.
\textit{Evidence: }[Research session 21, artifact 3]
\subsection{C10 -- Unresolved}
Neither a predictive improvement nor a biological mechanism can be claimed without biological-unit-aware baselines, feasible pre-outcome measurements, quantified uncertainty, and genuinely independent evidence.
The specified comparisons include stage-only, RIC/ECM-only, and known-signaling baselines with biological-unit splits, clutch-aware uncertainty, and negative controls. A destructive RIC assay cannot silently become an individual predictor for another animal followed to regeneration; a validated live surrogate or an explicitly narrower target is needed. A nonsignificant rescue is inconclusive without valid manipulation checks and enough precision to exclude a prespecified meaningful effect. Shared-study assays and reused publications are not independent biological validation.
\textit{Evidence: }[Project overview], [Research session 21, artifact 1], [Research session 21, artifact 3], [Research session 21, artifact 4]
\subsection{C11 -- Unresolved}
A stage-40 versus stage-46 comparison cannot establish why competence subsequently returns, and no verified later-competent Xenopus laevis condition supports the candidate mechanism.
The saved obligations require independent confirmation and sampling of a later competent condition in Xenopus laevis. Stage patterns or nutritional findings attributed to Xenopus tropicalis are not transferred as supporting evidence. The absence of this comparison limits the scope of any proposed loss-of-competence mechanism.
\textit{Evidence: }[Research session 21, artifact 1], [Research session 21, artifact 4], [Research session 21, artifact 5]
\subsection{C12 -- Limitation}
The available evidence supports an analytical identifiability result and untested mechanistic proposals, but neither a new biological mechanism nor a reproducible computational predictor has been demonstrated.
The evidence explicitly reports no expression reanalysis, fitted baseline, numerical biological effect, independent molecular result, or functional regeneration validation. Original-source applicability and novelty remain unassessed. Missing evidence is not a negative biological result, and model review is not independent experimental or formal verification.
\textit{Evidence: }[Research session 21, artifact 5], [Research session 21]
\section{Limitations}
\begin{itemize}
\item The central analytical result is conditional on observing only unlinked stage/time marginals without additional identifying restrictions. Neither the presence nor the absence of the required linkage has been verified in the actual GEO samples. This limits application of C3 and C4, not their constructed counterexamples.
\item The literature boundary is provisional. The original cited publications and later work have not been inspected sufficiently to call the receptivity-window hypothesis previously unreported; C1 therefore precludes a novelty claim.
\item Selective basal-epidermal ERK manipulation, matched exposure, live state measurements, and a timing schedule that separates wound age from response latency remain unvalidated capabilities. The model contrast in C6 does not establish experimental feasibility or a regrowth effect.
\item No empirical effect estimate, uncertainty interval, fitted baseline, independent validation, or verified later-competent Xenopus laevis comparison is available. The present outcome does not meet either the mechanism-discovery criterion or the lower-tier reproducible computational-result criterion.
\end{itemize}
\section{Open obligations}
\begin{itemize}
\item For C2, inspect original GEO series and sample records and build a row-per-library manifest covering species, developmental stage, collection time, animal or pool identity, biological versus technical replication, batches, units, assay availability, linkage, and regeneration labels. Restrict analysis to identifiable population contrasts if individual temporal linkage is absent.
\item For C1, complete an inspected, dated prior-art ledger through the investigation cutoff. Record species, experimental context, the interaction actually tested, and precise source locations; abandon the novelty claim if the proposed dependency is already established.
\item For C2 and C10, verify stage\textendash{}batch separability, freeze molecular modules and Xenopus laevis homeolog mapping, and distinguish changing cell abundance from within-cell state changes before fitting archival contrasts.
\item For C5\textendash{}C7, validate epidermal reporters and selective perturbations, operationally define receptivity and RIC-state entry, and establish a molecular sampling schedule separating wound age from post-pulse latency. Stop interpreting a timing contrast if exposure matching or selective delivery fails.
\item For C8, reformulate any persistence-specific redox test so that its temporal predictor can vary independently of the controlled exposure integral. For C9, specify and assess a spatial null that addresses anatomical gradients, autocorrelation, registration, and patch structure.
\item For C10, preregister prediction targets, feasible pre-outcome baselines, biological-unit splits, negative controls, primary metrics, multiplicity handling, and uncertainty estimation. Use a separate pilot to justify meaningful-effect bounds and clutch-aware sample size rather than treating a proposed margin or animal count as adequate.
\item For C5 and C10, obtain independent randomized molecular and functional regeneration observations with sham, vehicle, wound-size, closure, viability, and composition controls. Compare early-signal perturbation with a validated downstream bypass and reciprocal blockade where feasible; do not interpret rescue alone as identification of the decoder.
\item For C11, independently verify a later competent Xenopus laevis condition and test whether the proposed temporal relationship returns with competence. Do not import a developmental trajectory from Xenopus tropicalis.
\end{itemize}
\section{Frozen obligation reconciliation}
Unlisted obligations are not presumed resolved. The archived research snapshot retains complete source references and the full reconciliation.
\begin{itemize}
\item ChatGPT plan or app usage limit reached. Manage usage at https://chatgpt.com/settings/usage; no API-key fallback occurred. (HTTP 400) No strategy was saved. \textit{ (Not listed by author; not presumed resolved; 1 saved source reference(s))}
\item No research artifact was saved for this strategy. \textit{ (Not listed by author; not presumed resolved; 3 saved source reference(s))}
\item Acquire original GEO sample metadata and files; verify species, exact stages/times, independent animals versus pools, pairing, batches, replicate counts, outcome labels and assay availability before analysis. If only unlinked snapshots exist, stop the animal-level temporal-coupling analysis. \textit{ (Not listed by author; not presumed resolved; 3 saved source reference(s))}
\item Inspect the cited original papers and later literature through the investigation date; record exact evidence for or against prior reporting of the pulse-timing interaction. If already established, abandon its novelty claim. \textit{ (Not listed by author; not presumed resolved; 3 saved source reference(s))}
\item Validate basal-epidermal exposure and response measurements, noninjurious pulse delivery and stable response-score normalization. If early and late exposures cannot be matched, stop rather than interpret their contrast as a timing effect. \textit{ (Not listed by author; not presumed resolved; 3 saved source reference(s))}
\item Use a separate pilot to justify the minimum-interest margin, estimate between-animal and between-clutch variance, and choose a finite confirmatory sample size for the interaction. Do not select the response module or effect threshold using confirmatory outcomes. \textit{ (Not listed by author; not presumed resolved; 3 saved source reference(s))}
\item Obtain independent randomized molecular and regeneration results with uncertainty. Resolve cell composition, developmental progression, metabolic state and injury-severity alternatives; a pulse rescue alone does not identify the decoder. \textit{ (Not listed by author; not presumed resolved; 3 saved source reference(s))}
\item Animal-level prediction and any claim explaining later regained competence require additional longitudinal measurements and a verified later-competent X. laevis condition, respectively. \textit{ (Not listed by author; not presumed resolved; 3 saved source reference(s))}
\item Audit original GEO metadata, biological units, pairing, pools, batches, phenotypes and file availability. \textit{ (Not listed by author; not presumed resolved; 4 saved source reference(s))}
\item Inspect original and subsequent literature for exact overlap with the pulse-timing interaction. \textit{ (Not listed by author; not presumed resolved; 4 saved source reference(s))}
\item Validate pulse delivery, local exposure, basal-epidermal assays and normalization; separate wound age from post-pulse response latency. \textit{ (Not listed by author; not presumed resolved; 4 saved source reference(s))}
\item Justify the minimum-interest margin and determine confirmatory sample size using separate pilot animals and clutch-level variability. \textit{ (Not listed by author; not presumed resolved; 4 saved source reference(s))}
\item Obtain independent randomized molecular and regeneration results, with uncertainty and composition, metabolic, developmental and injury controls. \textit{ (Not listed by author; not presumed resolved; 4 saved source reference(s))}
\item Require individual measurement\textendash{}outcome links for animal-level prediction and a verified later-competent condition for regained-competence claims. \textit{ (Not listed by author; not presumed resolved; 4 saved source reference(s))}
\item Material and recoverable: the fixed 6-hour endpoint confounds pulse timing with elapsed time since treatment. Stage-dependent response kinetics can produce the proposed interaction without the claimed wound-age-dependent gate. \textit{ (Not listed by author; not presumed resolved; 2 saved source reference(s))}
\item Material and recoverable: neither accession structure nor prior reporting of the timing interaction has been inspected. The mathematical limitation cannot yet be attributed to the actual GEO samples, and novelty cannot be claimed. \textit{ (Not listed by author; not presumed resolved; 2 saved source reference(s))}
\item Material and recoverable: no independent molecular or regeneration measurements support the candidate, and no verified later-competent X. laevis condition tests regained competence. These are missing evidence, not refutations. \textit{ (Not listed by author; not presumed resolved; 2 saved source reference(s))}
\item The proposed timing contrast confounds wound age with response latency. \textit{ (Not listed by author; not presumed resolved; 2 saved source reference(s))}
\end{itemize}
111 further distinct obligations are retained in the archived reconciliation.
\section{Conclusions}
Unlinked early and late population measurements cannot generally establish whether individual animals with stronger early injury responses later activate more regenerative epidermal output. This conditional nonidentifiability result does not show that biological coupling is absent, and its application to the named datasets awaits the sampling audit.

The restricted receptivity model provides a specific counterfactual: matched early and late epidermal ERK supplementation has opposite transition-drive effects under early signal insufficiency and delayed receptivity. This is a model-level prediction whose biological interpretation depends on selective manipulation and separation of response latency from wound age.

The mechanism responsible for temporary loss and later recovery of Xenopus laevis regeneration competence remains unresolved. Discovery claims require source-checked novelty, numerical effects with uncertainty and appropriate baselines, and independent biological confirmation, including a verified later competent condition.

\textit{Claims discussed: }C2, C3, C5, C6, C7, C10, C11, C12
\appendix
\section{Selected source-linked research excerpts}
The following are bounded, selected excerpts of original research records. They are not exhaustive; consult the full frozen snapshot archived with this paper. They do not constitute independent verification.
\subsection{Research excerpt 1}
\textit{Source and field: }[Research session 21, artifact 4]: skeptic\_output/reasoning\_assessments\par
[\{"basis\_artifact\_ids": [], "disposition": "accepted", "objection": "", "rationale": "The boundary matches the project-supplied description and explicitly distinguishes that description from inspected literature. It does not claim verified novelty or independent replication.", "step\_number": 1\}, \{"basis\_artifact\_ids": [], "disposition": "accepted", "objection": "", "rationale": "The equations define coherent restricted candidate models.

They are not inferred from data, and the argument correctly recognizes that broad RIC-induction failure contains both alternatives.", "step\_number": 2\}, \{"basis\_artifact\_ids": [], "disposition": "accepted", "objection": "", "rationale": "With C and the untreated trajectory fixed, linearity of the integral gives the stated contrast. Constant c on only one pulse window yields exactly -kcQ or +kcQ.

Transfer to regrowth is explicitly conditional on the listed monotone, nonsaturated response assumption; it does not follow from the R and M differential equations alone.", "step\_number": 3\}, \{"basis\_artifact\_ids": [], "disposition": "accepted", "objection": "", "rationale": "This is a feasibility-gated intervention proposal, not an assertion that a validated X. laevis reagent exists.

Measuring activity rather than equating illumination durations appropriately addresses stage-dependent actuator response.", "step\_number": 4\}, \{"basis\_artifact\_ids": [], "disposition": "accepted", "objection": "", "rationale": "Under early-supported stage-40 receptivity, suppressing relevant early activity reduces cumulative transition drive whereas late suppression\ldots{}

\subsection{Research excerpt 2}
\textit{Source and field: }[Research session 21]: outcome/solution\_review/strategy\_findings/1/reasoning\_steps\par
[\{"argument": "Let E\_i denote an animal's early wound-response measurement, R\_i its later epidermal RIC/ECM-transition measurement, and S\_i developmental stage. An animal-level coupling estimand is the within-stage slope Cov(E,R|S)/Var(E|S), when the variance is positive. Measurements from different animals at different times estimate marginal distributions, not this joint distribution. Whether the actual GEO samples provide animal linkage, independent animals, pools, or only technical/library replication remains unverified.

Pool-level results must not be represented as animal-level prediction.", "basis\_artifact\_ids": [], "epistemic\_status": "deduction", "number": 1, "research\_artifact\_id": "[Research session 21, artifact 2]", "review": \{"basis\_artifact\_ids": [], "disposition": "accepted", "objection": "", "rationale": "Cov(E,R|S)/Var(E|S) is a defined within-stage linear-projection slope when the relevant moments exist and variance is positive. Unlinked sampling supplies marginals rather than animal correspondence. The step explicitly leaves the GEO sampling premise unverified."\}, "title": "Define the estimand and the unresolved sampling prerequisite"\}, \{"argument": "Within one stage, let U be uniform on [0,1].

Model I assigns E=U and R=U; Model II assigns E=U and R=1-U. Both models give exactly the same uniform marginal distributions for E and R. However, Var(E)=1/12, while Cov(E,R)=1/12 in Model I and -1/12 in Model II; the coupling slopes are therefore +1 and -1. Independently sampled early and late animals have identical observable distributions under these two\ldots{}

\subsection{Research excerpt 3}
\textit{Source and field: }[Research session 21, artifact 1]: skeptic\_output/reasoning\_assessments\par
[\{"basis\_artifact\_ids": [], "disposition": "accepted", "objection": "", "rationale": "The step explicitly treats the supplied biological background as a project constraint and distinguishes reference pointers from inspected sources. This accepts a stated premise and provenance boundary, not the underlying literature claims as independently established facts.", "step\_number": 1\}, \{"basis\_artifact\_ids": [], "disposition": "accepted", "objection": "", "rationale": "The proposed regression slope depends on a joint covariance, whereas separate marginal samples do not generally determine that covariance. Matching group labels does not create individual pairing.

The slope is meaningful where its requisite moments exist and early-response variance is positive.", "step\_number": 2\}, \{"basis\_artifact\_ids": [], "disposition": "accepted", "objection": "", "rationale": "For symmetric ±1 early responses, variance is 1. Setting the later response equal to or opposite the early response gives covariance and slope +1 or -1 while preserving both marginals. The two stated stage contrasts are therefore exactly 0 and -2. This is an analytic construction, not computation or frog data.", "step\_number": 3\}, \{"basis\_artifact\_ids": [], "disposition": "accepted", "objection": "", "rationale": "Stage-specific shifts and positive rescaling leave alternative pairings possible.

Group-mean correlations consequently do not recover within-stage individual covariance. Cell-state annotation alone likewise supplies no missing longitudinal identity or refractory sampling.", "step\_number": 4\},\ldots{}

\subsection{Research excerpt 4}
\textit{Source and field: }[Research session 21, artifact 2]: investigator\_output/reasoning\_steps\par
[\{"argument": "Let E\_i denote an animal's early wound-response measurement, R\_i its later epidermal RIC/ECM-transition measurement, and S\_i developmental stage. An animal-level coupling estimand is the within-stage slope Cov(E,R|S)/Var(E|S), when the variance is positive. Measurements from different animals at different times estimate marginal distributions, not this joint distribution. Whether the actual GEO samples provide animal linkage, independent animals, pools, or only technical/library replication remains unverified.

Pool-level results must not be represented as animal-level prediction.", "basis\_artifact\_ids": [], "epistemic\_status": "deduction", "title": "Define the estimand and the unresolved sampling prerequisite"\}, \{"argument": "Within one stage, let U be uniform on [0,1]. Model I assigns E=U and R=U; Model II assigns E=U and R=1-U. Both models give exactly the same uniform marginal distributions for E and R. However, Var(E)=1/12, while Cov(E,R)=1/12 in Model I and -1/12 in Model II; the coupling slopes are therefore +1 and -1. Independently sampled early and late animals have identical observable distributions under these two models. Thus no analysis of those unlinked measurements alone can distinguish even the sign of individual coupling.

Additional unlinked time points do not supply the missing individual correspondence without extra, testable structural assumptions. This construction is an identifiability proof, not a biological model of tadpole expression.", "basis\_artifact\_ids": [], "epistemic\_status": "deduction", "title": "Prove nonidentifiability from\ldots{}

\subsection{Research excerpt 5}
\textit{Source and field: }[Research session 21, artifact 3]: skeptic\_output/reasoning\_assessments\par
[\{"basis\_artifact\_ids": [], "disposition": "unresolved", "objection": "This entry is designated an assumption, while claim\_is\_conditional is false. It cannot be accepted as an established premise. Its accession descriptions and prior-art attributions remain background awaiting inspection.", "rationale": "The evidence boundary is appropriately disclosed, and the project-supplied references are correctly treated as provisional rather than inspected sources.", "step\_number": 1\}, \{"basis\_artifact\_ids": [], "disposition": "accepted", "objection": "", "rationale": "Direct arithmetic gives mean(I)=mean(M)=1/2, mean(I*M\_A)=1/2, and mean(I*M\_B)=0, hence covariances 1/4 and -1/4.

The conclusion is correctly restricted to non-identification by the specified marginals.", "step\_number": 2\}, \{"basis\_artifact\_ids": [], "disposition": "accepted", "objection": "Acceptance concerns the definition only. Biological inference still requires a defensible exchangeability model and sensitivity to spatial autocorrelation, registration error, and stratum choice.", "rationale": "T minus its specified null expectation is a coherent definition of excess alignment.

The text explicitly limits interpretation to the chosen constraints and does not claim reporter validity or control of all anatomy.", "step\_number": 3\}, \{"basis\_artifact\_ids": [], "disposition": "unresolved", "objection": "Neither the positive association nor its developmental loss-and-return pattern follows from the illustrative construction. They require biological observations; adjustment alone cannot establish a coupling defect.",\ldots{}

\subsection{Research excerpt 6}
\textit{Source and field: }[Research session 21, artifact 5]: skeptic\_output/reasoning\_assessments\par
[\{"basis\_artifact\_ids": [], "disposition": "accepted", "objection": "", "rationale": "The supplied context contains accession references but no original sample tables, matrices or publications. Treating sample structure and pairing as unverified follows from the available evidence.", "step\_number": 1\}, \{"basis\_artifact\_ids": [], "disposition": "accepted", "objection": "", "rationale": "The centered variables take values ±1 with equal probability. Same-sign pairing gives covariance +1, opposite-sign pairing gives -1, and early-response variance is 1 in both cases.

Swapping the assignments between stages reverses the slope interaction while preserving each marginal.", "step\_number": 2\}, \{"basis\_artifact\_ids": [], "disposition": "accepted", "objection": "", "rationale": "Observational equivalence proves nonidentifiability in the unrestricted marginal-only model. Increasing unlinked sample counts does not recover the missing joint distribution.

The stated qualifications concerning longitudinal links, structural restrictions and pool-level versus animal-level inference are appropriate.", "step\_number": 3\}, \{"basis\_artifact\_ids": [], "disposition": "accepted", "objection": "", "rationale": "The ledger accurately distinguishes project-supplied descriptions from inspected literature and makes no priority claim. It is a provisional source list, not a completed novelty audit.", "step\_number": 4\}, \{"basis\_artifact\_ids": [], "disposition": "challenged", "objection": "At 6 hours, early and late pulses have different response latencies. Transient induction or score decay can generate\ldots{}

\subsection{Research excerpt 7}
\textit{Source and field: }[Research session 8]: steps/0/provider\_input\_audit/request\_body/text/format/schema/properties/reasoning\_steps\par
\{"items": \{"\$ref": "\#/\$defs/ResearchReasoningStep"\}, "maxItems": 24, "title": "Reasoning Steps", "type": "array"\}

\subsection{Research excerpt 8}
\textit{Source and field: }[Research session 21]: steps/9/output/reasoning\_assessments\par
[\{"basis\_artifact\_ids": [], "disposition": "accepted", "objection": "", "rationale": "The boundary matches the project-supplied description and explicitly distinguishes that description from inspected literature. It does not claim verified novelty or independent replication.", "step\_number": 1\}, \{"basis\_artifact\_ids": [], "disposition": "accepted", "objection": "", "rationale": "The equations define coherent restricted candidate models.

They are not inferred from data, and the argument correctly recognizes that broad RIC-induction failure contains both alternatives.", "step\_number": 2\}, \{"basis\_artifact\_ids": [], "disposition": "accepted", "objection": "", "rationale": "With C and the untreated trajectory fixed, linearity of the integral gives the stated contrast. Constant c on only one pulse window yields exactly -kcQ or +kcQ.

Transfer to regrowth is explicitly conditional on the listed monotone, nonsaturated response assumption; it does not follow from the R and M differential equations alone.", "step\_number": 3\}, \{"basis\_artifact\_ids": [], "disposition": "accepted", "objection": "", "rationale": "This is a feasibility-gated intervention proposal, not an assertion that a validated X. laevis reagent exists.

Measuring activity rather than equating illumination durations appropriately addresses stage-dependent actuator response.", "step\_number": 4\}, \{"basis\_artifact\_ids": [], "disposition": "accepted", "objection": "", "rationale": "Under early-supported stage-40 receptivity, suppressing relevant early activity reduces cumulative transition drive whereas late suppression\ldots{}

\subsection{Research excerpt 9}
\textit{Source and field: }[Research session 21, artifact 2]: objective\_review/argument\_steps\par
[\{"argument": "Let E\_i denote an animal's early wound-response measurement, R\_i its later epidermal RIC/ECM-transition measurement, and S\_i developmental stage. An animal-level coupling estimand is the within-stage slope Cov(E,R|S)/Var(E|S), when the variance is positive. Measurements from different animals at different times estimate marginal distributions, not this joint distribution. Whether the actual GEO samples provide animal linkage, independent animals, pools, or only technical/library replication remains unverified.

Pool-level results must not be represented as animal-level prediction.", "basis\_artifact\_ids": [], "epistemic\_status": "deduction", "number": 1, "research\_artifact\_id": "[Research session 21, artifact 2]", "review": \{"basis\_artifact\_ids": [], "disposition": "accepted", "objection": "", "rationale": "Cov(E,R|S)/Var(E|S) is a defined within-stage linear-projection slope when the relevant moments exist and variance is positive. Unlinked sampling supplies marginals rather than animal correspondence. The step explicitly leaves the GEO sampling premise unverified."\}, "title": "Define the estimand and the unresolved sampling prerequisite"\}, \{"argument": "Within one stage, let U be uniform on [0,1].

Model I assigns E=U and R=U; Model II assigns E=U and R=1-U. Both models give exactly the same uniform marginal distributions for E and R. However, Var(E)=1/12, while Cov(E,R)=1/12 in Model I and -1/12 in Model II; the coupling slopes are therefore +1 and -1. Independently sampled early and late animals have identical observable distributions under these two\ldots{}

6 further research records are retained in the archived source snapshot.
\section*{Evidence and authorship note}
Snapshot SHA-256: 3236a8c66a96d713f28726271169e67c88c8a768f398ab3b9b504fee3cdcc74e
Generated by an AI author from the cited saved research snapshot. This document is not independent verification. No external bibliography was added.
\section*{Validation warnings}
\begin{itemize}
\item Evidence indexing omitted 3370 material saved research field fragment(s). Those fragments were not reviewed for this paper; the interpretation is provisional where the omitted evidence may matter. Consult the frozen source snapshot and omission ledger for full details.
\item The evidence index is a bounded selection with omitted notes and fragments. This plan does not imply exhaustive review of the saved research or the literature.
\item No original sample manifest, expression matrix, inspected full-text publication, or independent biological dataset is supplied here. Accession coverage and detailed prior-publication conclusions must not be upgraded from provisional descriptions to verified observations.
\item C3, C4, and C6 are labeled established only as conditional analytical deductions supported by explicit constructions or algebra. Agent review is not independent formal verification, experimental validation, or evidence that their assumptions describe the biological system.
\item The numerical values in the nonidentifiability example and the signed pulse-drive contrasts are constructed quantities, not gene-expression measurements, regeneration effects, or confidence bounds.
\item The saved timing proposal contains an unresolved latency confound. A contrast at one common wound age cannot alone establish a wound-age-dependent gate.
\item The saved redox persistence design contains a dependent-variable problem: matching an integral also fixes any score exactly proportional to that integral. Its proposed control therefore cannot establish a persistence-specific effect.
\item The spatial-alignment proposal lacks a validated structure-preserving null and linked individual outcome baseline. Marginal abundance differences or simple spatial permutations cannot establish its mechanism.
\item No biological null result is established. Missing measurements, incomplete analysis, and unsuccessful execution are not evidence against any candidate mechanism.
\item The evidence index explicitly omitted 3370 saved field fragment(s); see the generation coverage ledger and review the complete archived snapshot.
\item Structured paper indexing selected 466 source-linked text fragments from 15 archived sources without AI extraction calls. The complete frozen snapshot remains archived; this bounded index does not constitute review of every saved field.
\item C3 is labeled established, but cited saved evidence includes unverified, incomplete, or conjectural material.
\item C3 is labeled established while the frozen evidence contains counterexample or contradiction language; inspect the cited sources before accepting it.
\item C4 is labeled established, but cited saved evidence includes unverified, incomplete, or conjectural material.
\item C4 is labeled established while the frozen evidence contains counterexample or contradiction language; inspect the cited sources before accepting it.
\item C6 is labeled established, but cited saved evidence includes unverified, incomplete, or conjectural material.
\item C6 is labeled established while the frozen evidence contains counterexample or contradiction language; inspect the cited sources before accepting it.
\item The frozen evidence contains open proof/research obligation(s) requiring reconciliation: 121 item(s) are not listed by the author and are not presumed resolved; 8 author-listed item(s) have no exact frozen-record match. See the obligation reconciliation.
\item A diagram's graph could not be laid out without likely crossings or overlap; its complete nodes, boundaries, and relationships are shown in a readable text representation.
\end{itemize}
\section*{Typesetting verification note}
Needs review: TeX verification is incomplete. The authored research content and safe accepted typesetting corrections were retained. Review the saved proofreading diagnostics before relying on the layout.
\end{document}
